\documentclass[10pt]{article}
\usepackage[a4paper,margin=25mm]{geometry}
\usepackage{amsmath,amssymb,amsthm,mathrsfs,longtable,hyperref}
\emergencystretch=2em
\newtheorem{theorem}{Theorem}[section]
\newtheorem{lemma}[theorem]{Lemma}
\newtheorem{proposition}[theorem]{Proposition}
\newtheorem{corollary}[theorem]{Corollary}
\newcommand{\PSH}{\operatorname{PSH}}
\newcommand{\FS}{\operatorname{FS}}
\newcommand{\vol}{\operatorname{vol}}
\newcommand{\MA}{\operatorname{MA}}
\newcommand{\an}{\mathrm{an}}
\newcommand{\OO}{\mathcal O}
\newcommand{\E}{\mathcal E}
\title{Original-carrier foundations and trace estimates in big rational classes\\
\large An incomplete research paper for the unchanged original P16.5}
\author{Research reconstruction from the two explicit mathematical grants}
\date{6 October 2026}
\begin{document}
\maketitle
\begin{abstract}
We reconstruct a relative-floor decreasing-closure argument using the granted
relative multiplier-ideal net, before constructing the intrinsic section-volume
energy. Complete moving systems recover every power of the ordinary positive
product, including every fixed signed Cartier test. An original-section trace
inequality retains its missing-rank term. It transfers finite nef-model
contraction to the entire genuine finite-energy cone and gives an intrinsic
relative integrability criterion and coercivity. A separate outer derivative
construction yields target-independent positive measures of volume mass.
These are partial research deductions, not a completed finite-energy
Monge--Amp\`ere equation theorem. Unbounded weighted identification, arbitrary
target attainment and comparison remain unproved in this paper. The additional
zero-volume equation range proved here is exactly the zero numerical class.
No zero mass is divided by volume or normalized to a probability.
\end{abstract}

\section{Scope, premises and the incomplete canonical result}
Let $k$ be any algebraically closed field of characteristic zero, with its
trivial absolute value. Let $X$ be any smooth integral projective $k$-variety,
$n=\dim X\geq1$, $K=X^\an$, and $L$ an actual big rational line bundle.
Set $\theta=c_1(L)$ and $V=\vol_X(L)>0$.
There is no restriction on cardinality, dimension, section-ring finite
generation, invariance or measure support. An actual rational bundle means
an element of $\operatorname{Pic}(X)\otimes\mathbb Q$; tensor powers are
cleared whenever sections occur. Numerical equality never replaces an actual
tensor identity. Scalar potentials are measured from the canonical trivial
metric. The generic trivial valuation is $v_0$.

Only two mathematical bodies are task-granted. G0 is the complete supplied
P16.1 body in its actual ample rational-bundle domain, over these fields and
reduced irreducible geometrically unibranch projective carriers. Its genuine
ample regularization, normalized compactness, all-point order determination,
raw intersection operator, energy theory and equation may be used in that
domain. G0 grants no equation for $L$. G1 is the complete supplied relative
asymptotic-multiplier-ideal generation theorem on smooth equivariant resolved
modifications of $X\times\mathbb P^1$, with actual vertical rational Cartier
data and an actual linearized ample twist. Its single full generically finite
FS net is granted; its stages may have supported poles. G1 grants no general
closure, envelope identity, energy, comparison or equation. All other supplied
papers are unreviewed data. Their arguments are reconstructed below where used.

We use ordinary algebraic characteristic-zero background: resolution and
principalization, extraction of a divisorial Gauss valuation on an equivariant
model, Kodaira's ample-plus-effective decomposition, ample global generation,
proper birational normal section descent, section-volume asymptotics, Fujita
approximation, and the positive-product construction and volume differential
of Boucksom--Favre--Jonsson (BFJ). For a smooth projective $d$-fold,
\[
 D\vol_\alpha(H)=d\langle\alpha^{d-1}\rangle\cdot H
\]
on its big cone. The positive products take values in the weak Weil numerical
space; their pseudoeffective order can be paired with nef Cartier classes.
The ordinary nef-product and volume-root concavity results are used in their
algebraic domains. They are not nonample NA equation grants. The precise
primary read scope and underlying published background are recorded at the end.

\paragraph{Canonical result of this submission.}
The all-point cone, minimal reference and scalar finite-energy construction in
Sections 2--6, and the cofinality, trace, contraction and intrinsic-dual
deductions in Sections 7--10, are partial proof drafts relative to the two
grants and the explicitly scoped ordinary background. Section 11 constructs
an outer derivative measure on every finite-energy endpoint. It has not been
identified here through every unbounded weighted test with the required
genuine finite-energy Monge--Amp\`ere operator. We therefore claim neither the
full big-class equation nor original P16.5. The canonical result has
\texttt{candidate=false}; Section 13 states exact remaining assertions.
The original task, complete family, grants and all supplied candidate bodies
are preserved literally in the accompanying evidence, not paraphrased away.

\section{The genuine numerical cone and elementary actual-bundle operations}
A full generically finite FS function on an actual rational bundle $B$ is
\[
 \phi=\frac1m\max_j\{\log|s_j|+\lambda_j\},\quad
 0\ne s_j\in H^0(X,mB),\quad\lambda_j\in\mathbb R.
\]
Sections need not generate $mB$. Such a function is continuous into
$\mathbb R\cup\{-\infty\}$ and has finite generic value. Define
$\PSH(\theta)$ to consist of the all-point decreasing limits of one net of
these actual functions on bundles $B_i$ with $c_1(B_i)\to\theta$, whose
limiting generic value is finite. This is a definition by actual functions on
$K$, not a completion and not a twice-closed family. Weak convergence means
convergence at rational divisorial valuations, together with $v_0$.

The maximum of the canonical norms of a generating system of an actual
ample power is one everywhere: at the center of each analytic point one
generator is a unit. Tensoring with this weight-zero system embeds an FS
function on $B$ into $B+C$, for actual ample $C$, without changing its scalar.
If $c_1(B_i)\to\theta$ and $\Omega-L$ is ample, a tail satisfies
$\Omega-B_i$ ample. The same scalar net therefore embeds into the one
actual ample cone $\PSH(\Omega)$. If finite real-valued stages are required
there, append generators at weight $-mN$ and use the product net
$\max(\phi_i,-N)$. This is one decreasing net; it is not an assertion of
nonample decreasing closure.

Finite maxima and convex combinations also have single defining nets.
Tensor two incoming section systems into one actual target whose ample
differences tend numerically to zero, and take their union for a maximum.
For a real coefficient $0<r<1$, put both incoming functions below a common
finite $C$, choose positive rational $r_j\uparrow r$ and
$s_j\uparrow1-r$, and use
\[
 r_j(\phi_i-C)+s_j(\psi_l-C)+C.
\]
These functions decrease on the product index, are FS after clearing powers,
and have the desired all-point limit, including poles. No pole is subtracted
from another pole. Constants and rational positive scaling are immediate.
Generic normalization is valid because canonical section norms are at most
one, so every FS stage and its endpoint have supremum equal to their value
at $v_0$. G0 in the common actual ample class gives finite divisorial values
and all-point order determination. Its normalized compact space gives, for
every fixed divisorial $z$, a finite uniform bound
\[
 -T_\Omega(z)\leq u(z)\leq0
 \quad(u\in\PSH(\theta),\ u(v_0)=0).                 \tag{2.1}
\]
No equality of ample and $\theta$ energies is asserted.

\section{The G1 net identifies model envelopes before closure}
Fix an actual big rational $B$ and a rational continuous model obstacle $f$.
Represent $f$ on a smooth equivariant resolved model with generic bundle $B$,
vertical Cartier divisor $D$, and full central fiber $F=\sum b_EE$.
G1 supplies complete relative base ideals $\mathfrak a_m$, their asymptotic
multiplier ideals $\mathfrak b_m$, one actual ample twist $\mathcal C$ with
bounded $\phi_{\mathcal C}\geq0$, and the one decreasing net
\[
 \Psi_m=f+\frac1m\phi_{\mathfrak b_m}
               +\frac1m\phi_{\mathcal C}\in\FS(B+C/m).
                                                               \tag{3.1}
\]
The indices with nonempty complete $B$ section spaces contain a cofinal
divisibility subset. Generation itself holds in every admissible degree.
Writing $Q_B(f)$ for the supremum of all full FS minorants on the actual $B$,
G1 gives
\[
 Q_B(f)\leq\Psi:=\inf_m\Psi_m\leq f,
 \qquad\Psi\in\PSH(c_1(B)).                           \tag{3.2}
\]
The last membership uses this one granted net, not decreasing closure.

We prove $\Psi=P_B(f)$, the all-point envelope of genuine competitors.
Write the actual Kodaira identity $B=C_0+E$ with $C_0$ ample and $E$ effective,
and $q_E=\log|s_E|\leq0$, after clearing powers. Given genuine $\tau\leq f$,
take its original defining net. A directed tail can be represented in
$B+\epsilon_i C_0$, $\epsilon_i>0$ rational tending to zero, by actual ample
difference tensoring. Such valid pairs are directed by making $\epsilon$
smaller and choosing an index in its ample-cone tail. For every fixed
$a>0$, compactness of $K$ makes a tail of its scalar stages $h_i$ satisfy
$h_i<f+a$ everywhere: the closed sets $\{h_i\geq f+a\}$ are nested with
empty intersection. The actual identity
\[
 (B+\epsilon_i C_0)+\epsilon_i E=(1+\epsilon_i)B
\]
therefore gives $\chi_i=(h_i+\epsilon_i q_E)/(1+\epsilon_i)\in\FS(B)$.
Put $M=\|f\|_\infty$. The nonnegative scalar
$a+\epsilon_i M$ is an upper bound for $\chi_i-f$; hence
$\chi_i-a-\epsilon_i M\leq Q_B(f)$ all-point. At each divisorial $z$,
$q_E(z)$ is finite and $\chi_i(z)\to\tau(z)$. Consequently
$\tau(z)\leq Q_B(f)(z)+a\leq\Psi(z)+a$. Let $a\downarrow0$.
Embed the already genuine $\tau,\Psi$ in one actual ample cone and apply
G0's order determination: $\tau\leq\Psi$ on every point of $K$.
Since $\Psi\leq f$ is itself genuine, this proves $P_B(f)=\Psi$.

In particular $V_B=P_B(0)$ is genuine, has generic value zero, and dominates
every normalized genuine $B$ potential. Normalize any sufficiently late
stage of (3.1) to $\psi\leq0$. Its finite generic shift gives
\[
 V_B\leq\psi+c_\psi                                \tag{3.3}
\]
for a finite $c_\psi$. These particular references exist before the next
section uses them. No full Addendum envelope theorem is an input here.

\section{Relative floors: one all-point defining net}
\begin{lemma}[Compact comparison across old poles]\label{polelemma}
Suppose an FS function $F$ obeys
$F\geq(1-\delta)\psi-k$, with $\psi\leq0$ FS, $0<\delta<1$.
If $h$ is usc and $h\leq\psi+C$, the function equal to $h-F$ where $F$
is finite and assigned $-\infty$ at the poles of $F$ is usc. For a decreasing
net of such $h$ with a common $C$, these comparison functions decrease.
If their infimum is strictly negative everywhere, a tail is negative everywhere.
\end{lemma}
\begin{proof}
Off the poles the difference is usc and at most $\delta\psi+C+k$.
A pole of $F$ forces $\psi=-\infty$, and thus $h=-\infty$ there. Extended
continuity of $\psi$ makes this upper bound tend to $-\infty$ at that point.
The assigned value is a comparison extension, not an arithmetic value for
$(-\infty)-(-\infty)$. The nonnegative superlevel sets of a decreasing
comparison net are nested compact sets. If their intersection is empty,
directedness and compactness make one of them, and then a tail, empty.
\end{proof}

\begin{theorem}[Genuine decreasing closure]\label{closure}
For a decreasing net $u_i\in\PSH(\theta)$ whose infimum $u$ has finite
generic value, $u$ has one actual all-point decreasing FS defining net with
classes tending to $\theta$.
\end{theorem}
\begin{proof}
Restrict to a tail with $u_i\leq M$. Fix an actual ample $A$, and write
$T_\eta=L+\eta A$ for rational $\eta>0$. Choose $0<\epsilon<\eta/4$,
$B=L+\epsilon A$, and a late normalized reference stage $\psi$ from (3.1)
on an actual bundle $R$, such that $T_\eta-R$ is ample. Choose rational
$\delta>0$ sufficiently small that $T_\eta-(1-\delta)R$ stays ample.
A special stage is an FS function $F$ on the specified $T_\eta$ with
\[
 F\geq u+\kappa,\qquad F\geq(1-\delta)\psi-k,       \tag{4.1}
\]
where $\kappa>0$ is rational and $k$ finite. It exists by taking a late
defining FS majorant $h$ of some $u_i$, tensoring it into $T_\eta$, tensoring
the floor into that same actual bundle, and taking their maximum.

For each fixed old reference $B,\psi$, all $u_i$ embed into $B$, and all
sufficiently late defining stages of any fixed $u_i$ do so too. Their generic
values are bounded by $M+1$. Minimality and (3.3) give the common bound
\[
 h\leq V_B+M+1\leq\psi+c_\psi+M+1.                \tag{4.2}
\]
This algebraic bound supplies the hypothesis of Lemma~\ref{polelemma}.

Take finitely many old special stages $F_l$, with their complete data and
guards $\kappa_l$. Choose new $\eta$ as small as desired, below all their
$\epsilon_l$, and a new normalized reference stage $\psi$ on $R$ sufficiently
close to $L$ that $B_l-R$ is ample for all $l$. Thus
$\psi\leq V_{B_l}\leq\psi_l+c_l$. Choose
$0<\delta<\min_l\delta_l$ also satisfying the new target amplitude condition.
For $b=(1-\delta)\psi-k$, where $F_l$ is finite,
\[
 b-F_l\leq(\delta_l-\delta)\psi_l+(1-\delta)c_l+k_l-k.
\]
Since $\psi_l\leq0$, one large $k$ gives $b\leq F_l$ for every $l$.
At a pole of $F_l$, the reference inequalities already force $b=-\infty$.

Choose $0<\kappa<\min_l\kappa_l/4$. Apply the compact comparison lemma to
$u_i+2\kappa$. At every finite old value its limiting difference is at most
$2\kappa-\kappa_l<0$; at an old pole the extension is $-\infty$.
One late $i$ satisfies all old strict comparisons. Apply the lemma again to
a late defining FS tail of this fixed $u_i$, using (4.2) separately for each
old reference. Choose one such stage $h\geq u_i$ with
$h+\kappa\leq F_l$ for all $l$, and with actual bundle close enough to $L$
to tensor into the new $T_\eta$. Then
\[
 F=\max(h+\kappa,b)\leq F_l\ (\forall l),\qquad F\geq u+\kappa.
                                                               \tag{4.3}
\]
The special stages, ordered by scalar decrease and $\eta$ decrease, form a
directed net. After any finite constraints its $\eta$ can be arbitrarily small.

Its infimum is $u$ all-point. Given $x$ and a finite $t>u(x)$, choose
$u_i(x)<t$ with a margin, choose the target gap $\eta$ before selecting a
sufficiently late defining stage $h$ fitting it, and choose the guard so
$h(x)+\kappa<t$. A large $k$ makes the floor below $t$ at $x$, or it is
already $-\infty$ there. This creates a special stage with $F(x)<t$.
Directedness subsequently imposes every old constraint. The same argument
works for $u(x)=-\infty$ with any finite $t$. Every stage stays above $u$.
Its numerical classes are $\theta+\eta a\to\theta$ and its generic endpoint
is the stipulated finite value. This is the single net required by the definition.
\end{proof}

An incoming upper stage is never multiplied by $1+\delta$; such scaling can
cross below the desired endpoint. The weakening $1-\delta$ applies solely
to the pole floor. The guard is a separate positive scalar.

\section{Envelopes and normalized compactness on the original carrier}
For usc bounded-above $f$ with a genuine minorant of finite generic value,
continuous majorants have all-point infimum $f$ on compact Hausdorff $K$.
Indeed a continuous cutoff supported in a neighborhood where $f<t$ gives a
majorant with value $t>f(x)$ at $x$; a global upper bound handles its exterior.
Rational model functions are a uniformly dense lattice in $C(K)$, as in the
actual ample G0 foundation. Upper uniform approximation and finite minima
give directed rational model majorants with the same infimum. Their envelopes
from Section 3 decrease and keep a finite generic lower bound. Theorem~\ref{closure}
makes their all-point limit genuine and below $f$. Every competitor lies below
each model envelope, so the limit is exactly $P_\theta(f)$. In particular the
usc supremum of any bounded-above nonempty genuine family is genuine: its
envelope dominates the raw supremum and hence its usc regularization; the
obstacle inequality gives equality.

For clarity we verify the ample nonnegligibility input to compactness rather
than importing big compactness. On a smooth SNC equivariant model with
central fiber $F=\sum b_jE_j$, let $x_j=b_j^{-1}\operatorname{ord}_{E_j}|_{k(X)}$.
For every FS function $h$ on any actual bundle and every original analytic
point $x$, its Gauss semivaluation has a center on this model. At that center,
$t$ is a unit times $\prod z_j^{b_j}$. Put
$\lambda_j=b_jw(z_j)\geq0$, so $\sum\lambda_j=w(t)=1$.
Factoring the pulled local section by its central prime orders gives
\[
 h(x)\leq\sum_j\lambda_jh(x_j).                     \tag{5.1}
\]
Its remaining local factor has nonnegative semivaluation, possibly infinity.
The maximum of weighted sums is at most the weighted sum of maxima. Passing
through the original defining net gives (5.1) for any genuine $h$, because the
finite vertex limits are finite. Supported semivaluations are included.
For a model function $q$, its value is exactly the same barycentric sum.
Consequently $\sup(h-q)=\max_j(h-q)(x_j)$.

Fix a rational divisorial $z$ and extract its Gauss valuation on such a model.
In one fixed actual ample cone, G0 compactness and finite divisorial values
give a finite uniform spread bound $C$ over this finite vertex set, unaffected
by constants. Choose a rational vertical model function $q$ with $q(z)=0$
and $q(x_j)>2C$ at all other vertices. Equation (5.1) makes
\[
 \sup_K(h-q)=(h-q)(z)                              \tag{5.2}
\]
for every genuine function in that ample cone. For a bounded-above family
$(h_i)$ it follows that $\sup h_i\leq\sup h_i(z)-q(z)+q$; continuity of $q$
preserves this bound under usc regularization. Evaluation at $z$ proves
$(\sup h_i)^*(z)=\sup h_i(z)$. This is the primary separator mechanism of
BJ Lemmas 2.12, 2.28--2.29, 4.26 and 5.7/Theorem 5.6, applied only after
actual ample embedding. Here (5.1) also supplies a direct SNC maximum proof.

Embed the normalized $\theta$ cone in one normalized actual ample cone.
For a net converging divisorially in the latter, regularized tail suprema are
genuine $\theta$ functions by the envelope property just proved. They decrease,
have generic value zero, and by ample nonnegligibility their divisorial infimum
is the prescribed coordinate limit. Theorem~\ref{closure} makes their
all-point infimum genuine $\theta$. Ample all-point determination identifies
it with the ample limit. Thus the normalized $\theta$ cone is a closed
subset of the compact ample cone, hence compact. This is a net argument;
there is no field-countability or sequential-compactness assumption.

\section{Section-volume energy and finite model derivative measures}
Let $\rho:Y\to X\times\mathbb P^1$ be a smooth equivariant modification,
product away from zero, $F=\sum b_jE_j=\rho^*\{0\}$, and $D$ vertical rational
Cartier. Write $f_D(x_j)=d_j/b_j$ and
$\alpha_D(c)=\rho_X^*L+D+cF$. Choose $aF\leq D\leq bF$ coefficientwise.
With $R_m=H^0(X,mL)$, $r_m=\dim R_m$, the actual relative section module is
the full graded lattice
\[
 t^{-ma}R_m[t]\subset\Lambda_m\subset t^{-mb}R_m[t],\qquad
 \Lambda_m=\bigoplus_jt^{-w_{mj}}k[t]s_{mj},\quad ma\leq w_{mj}\leq mb.
                                                               \tag{6.1}
\]
The inclusions follow from the actual effective divisor inequalities,
$\rho_*\OO_Y=\OO_{X\times\mathbb P^1}$ and projection formula. The flag in
$R_m$ furnished by the homogeneous lattice terms gives its adapted basis;
no finite generation of the full section ring is assumed. Gluing to the
product lattice at infinity gives $\bigoplus_j\OO(w_{mj})$ and hence
\[
 h^0(Y,m\alpha_D(c))=\sum_j\max(w_{mj}+mc+1,0).      \tag{6.2}
\]
For $c>-a$ this equals $(mc+1)r_m+\sum_jw_{mj}$. True full-section volume
asymptotics, first for fixed rational actual bundles and then continuity, give
\[
 \vol_Y(\alpha_D(c))=(n+1)cV+(n+1)\mathscr E(f_D)
 \quad\text{on one eventual half-line}.              \tag{6.3}
\]
The normalized weight sums therefore have an actual limit equal to this
intercept. It is birationally invariant, monotone, shifts by $cV$, and is
$V$-Lipschitz in the original model sup norm. Applying volume-root concavity
to (6.3), expanding at $c\to\infty$ with each fixed intercept, proves
ordinary concavity of $\mathscr E$.

For large $c$ put $q_j(D)=\langle\alpha_D(c)^n\rangle\cdot E_j$.
Effectivity of $E_j$ and the right volume differential show $q_j(D)\geq0$.
For $D+tE_j$ in a fixed small interval, (6.3) holds on a common half-line;
differentiation in $t$ proves that each coefficient, not merely their sum,
is independent of large $c$. Differentiation in the $F$ direction gives
\[
 \mu_D=\sum_jb_jq_j(D)\delta_{x_j}\geq0,
 \quad\mu_D(K)=\sum_jb_jq_j(D)=V,
 \quad D\mathscr E_{f_D}(f_H)=\int f_H\,d\mu_D.        \tag{6.4}
\]
If $p:Y'\to Y$ is a higher model and $G$ an effective exceptional prime,
normality gives $p_*\OO_{Y'}(rG)=\OO_Y$ for integral $r\geq0$. Sections and
volumes of $p^*\alpha+sG$ are unchanged for rational $s\geq0$. The volume
differential at the pulled big class has zero exceptional coefficient.
Strict transforms retain their old coefficient by pullback and projection.
Thus (6.4) is invariant under model refinement, with all $b_j$ retained.

For $f\in C(K)$ and $0\ne s\in R_m$ define
\[
 W_m(f;s)=\inf_K(mf-\log|s|).
\]
At a section pole the bracket is $+\infty$. The value is finite, the superlevel
conditions together with zero are linear subspaces, and they form a finite
closed-endpoint filtration. With an adapted basis set
$e_m(f)=n!m^{-(n+1)}\sum_j\lfloor W_m(f;s_{mj})\rfloor$.
The trace is intrinsic. For model $f_D$, normal section extension identifies
these floors exactly with (6.1). If $\|f-g\|_\infty\leq\epsilon$, flag
comparison gives
\[
 |e_m(f)-e_m(g)|\leq(n!r_m/m^n)(\epsilon+1/m).
\]
Uniform model density and $n!r_m/m^n\to V$ give a limit $\mathscr E(f)$ for
every original continuous $f$. This is the unique $V$-Lipschitz extension.
Each section term is a genuine $\theta$ minorant, so for every real $\lambda$
\[
 W_m(f;s)\geq\lambda\iff m^{-1}(\log|s|+\lambda)\leq P_\theta(f).
                                                               \tag{6.5}
\]
Therefore the energy is ordered by, and depends only on, its all-point envelope.

Define the actual relative energy and entire finite-energy space by
\[
 E(u)=\inf_{f\in C(K),\ f\geq u}\mathscr E(f),\qquad
 \E^1(\theta)=\{u\in\PSH(\theta):E(u)>-\infty\}.       \tag{6.6}
\]
This agrees with $\mathscr E(f)$ on $P_\theta(f)$ by (6.5). Its intrinsic
reference is $V_\theta=P_\theta(0)$; $E(V_\theta)=0$, $E(u+c)=E(u)+cV$.
If $u_i\downarrow u$, the compact closed sets $\{u_i\geq f+\epsilon\}$
are eventually empty for each continuous $f\geq u$. This gives
$\inf_iE(u_i)\leq\mathscr E(f)+\epsilon V$ and proves decreasing energy
continuity, also at $-\infty$. The genuine floors
$\max(u,V_\theta-N)$ decrease all-point to normalized $u$ and have energy in
$[-NV,0]$. Thus finite energy is exactly finite limiting relative-floor cost;
it does not require bounded difference from the reference.

Concavity and (6.4) give finite-coordinate global planes
\[
 E(u)\leq A_D(u):=\mathscr E(f_D)+\int(u-f_D)\,d\mu_D,
 \qquad E(u)=\inf_D A_D(u).                           \tag{6.7}
\]
For the first inequality pass decreasing continuous majorants at the fixed
finite divisorial support. For equality, take model majorants of $u$; their
nonpositive plane corrections make the infimum at most (6.6). The index $D$
is independent of $u$. Hence $E$ is concave and weakly usc, and normalized
sets $E\geq-B$ are compact sets of actual finite-energy functions.

\section{Complete moving systems and every fixed signed test}
Let $Z$ be smooth integral projective of dimension $d$, and let $N$ be an
actual big rational bundle. In an admissible nonempty degree $m$, resolve
the base ideal of the \emph{full} system, and write the actual identity
\[
 p_m^*(mN)=M_m+G_m,\qquad A_m=M_m/m.                 \tag{7.1}
\]
Here $M_m$ is globally generated and $G_m$ effective. Common dominations
are always used for comparisons. Multiplication of complete systems gives
$A_{m\ell}\geq A_m$ in the effective divisor order. A factorial divisibility
subnet clearing denominators is cofinal in every fixed admissible degree;
bigness supplies nonempty degrees on this subnet. No finite generation of
the section ring is involved.

We verify the positive-product cofinality, including lower powers. On one
fixed smooth determination let $\beta_1,\ldots,\beta_p$ be rational nef
competitors whose remainders in $N$ are pseudoeffective. Choose an actual
ample rational $Q_0$ with $N-Q_0$ big on that determination. For rational
$\epsilon>0$ put $Q_i=(1-\epsilon)\beta_i+\epsilon Q_0$. Choose actual
rational representatives for these classes. Each $Q_i$ is ample and
$N-Q_i$ is big: its class is the sum of $(1-\epsilon)(N-\beta_i)$ and the
big class $\epsilon(N-Q_0)$. In one sufficiently divisible degree $q$,
all $qQ_i$ are generated and all actual bundles $q(N-Q_i)$ have nonzero
sections with effective divisors $R_i$. Multiplication by these sections
injects their generated systems into the full system $|qN|$. Proper
birational normality identifies sections on this determination with the
original sections. Consequently $G_q\leq R_i$, and $A_q\geq Q_i$.

Intersections of nef bundles with an effective difference have the correct
monotonicity. Hence $A_q^p\geq Q_1\cdots Q_p$ in the weak effective-cycle
order. Let $\epsilon\downarrow0$ on this fixed determination. The least
upper bound of the moving products therefore equals the defining positive
product. The directed compactness argument of BFJ, Lemma 2.6 and Lemma 2.7,
converts this equality of upper bounds into weak convergence:
\[
 A_m^p\longrightarrow\langle N^p\rangle,
                 \qquad 1\leq p\leq d.             \tag{7.2}
\]
Real nef competitors follow by approximation inside the ample cone before
the preceding limit. This proof also handles several different nef
competitors; it is not merely a top-volume computation.

For a fixed signed Cartier test $H$ on a fixed determination, choose ample
rational $H_1,H_2$ there with $H=H_1-H_2$ numerically. Pair (7.2) with
$H_1$ and $H_2$ separately, then subtract. In particular
\[
 A_m^{d-1}\cdot H\longrightarrow
                  \langle N^{d-1}\rangle\cdot H.    \tag{7.3}
\]
No effective-order argument is applied directly to a signed $H$. The
test determination and the two ample tests are fixed before $m$ tends to
infinity. In (7.1), all equalities of bundles and systems are actual;
numerical classes enter only their intersection and positivity estimates.

\section{Horizontal fixed factors and the missing-rank trace term}
Fix a nonpositive rational model obstacle $f=f_D$, with $f(v_0)=0$, on a
smooth equivariant determination $\mathcal X$. Write
$F=\sum b_EE$ with the \emph{actual} central multiplicities. Let $w_{mj}$
be the full relative-lattice weights of Section 6, and choose one integer
$c$ larger than their uniform linear bound in absolute value. The absolute
system of $\alpha_D(c)$ has the actual basis
\[
 t^{-(w_{mj}+mc)+r}s_{mj},\qquad
                        0\leq r\leq w_{mj}+mc.      \tag{8.1}
\]
At infinity, the endpoint $r=w_{mj}+mc$ recovers every $s_{mj}$. At any
nonzero finite time their span also recovers the full generic system.
Thus this is an assertion about the entire lattice, including all its
weight spaces, rather than a selected graded subseries.

Resolve $|mL|$ on $X_m$ and write
\[
 p_m^*L=A_m+T_m,\quad mA_m\text{ generated},\quad mT_m\geq0.  \tag{8.2}
\]
On a common smooth equivariant determination over $X_m\times\mathbf P^1$,
every section in (8.1) has the actual factor $s_{mT_m}$. Remove that factor
\emph{before} principalizing the residual ideal. Because all sections of
$mA_m$ are then represented away from zero, the residual ideal is the unit
ideal there. Its resolved divisor $G_m^{\rm vert}$ is effective and is
supported at zero. The moving bundle is exactly
\[
 C_m=\rho_X^*A_m+D+cF-G_m^{\rm vert}/m,              \tag{8.3}
\]
and is globally generated, hence nef. Its bounded relative model
potential $\psi_m$ is in the actual generic nef bundle $A_m$. On this
determination the complete-system section function is
\[
 \phi_{m,D}=q_{T_m}+\psi_m,\qquad
                     \phi_{m,0}=q_{T_m}.            \tag{8.4}
\]
The logarithm $q_{T_m}$ is the canonical section norm. These are pulled
section identities. We neither descend $\psi_m$ arbitrarily nor assign
finite energy to the possibly polar physical function $\phi_{m,D}$.
Since all complete-system weights are nonpositive, $\psi_m\leq0$.
Its generic value $g_m$ tends to zero. To see the latter without an
envelope identity for complete systems, use the section minorants in
Section 3 to approximate the generic value of $P_\theta(f)=0$. Powers of
each such section belong to the full systems on all their multiples.

Apply Section 7 to both $L$ on $X$ and $\alpha_D(c)$ on $\mathcal X$:
\[
 V_m=A_m^n\to V,\quad C_m^{n+1}\to\vol(\alpha_D(c)),\quad
 C_m^n\cdot H\to\langle\alpha_D(c)^n\rangle\cdot H   \tag{8.5}
\]
for every fixed signed Cartier $H$. The relative nef energy is
\[
 e_m(\psi_m)=C_m^{n+1}/(n+1)-cV_m\longrightarrow E(P_\theta(f)).
                                                               \tag{8.6}
\]
This auxiliary bounded energy is an intersection quantity. Positivity and
its mixed-measure formulas follow by adding actual ample bundles on each
fixed smooth total model and taking the finite polynomial limit of G0.
Every central coefficient is $b_E$ times the appropriate restriction
intersection. This uses neither a general big equation nor an unbounded
nef equation. A sufficiently large time twist makes a relatively nef
model with nef generic bundle absolutely nef: horizontal curves are
controlled by $|D|\leq MF$, and vertical curves by relative nefness.
Ordinary nef volume equals its top intersection. The lattice computation
therefore identifies its section trace with this same energy.

Here is the comparison with an original continuous obstacle, with its
error retained. Fix $m$ and a bounded nonpositive relative nef model
potential $\psi$ on $A_m$. Suppose that an original $f\in C(K)$ dominates
the physical section function $q_{T_m}+\psi$ after pullback. In admissible
degree $\ell$, multiplication by $s_{\ell T_m}$ injects
\[
 H^0(X_m,\ell A_m)\hookrightarrow H^0(X,\ell L).    \tag{8.7}
\]
For every section in this subspace its original obstacle weight is at
least its auxiliary metric weight. This follows from analytic surjectivity
and the literal section identity in (8.7). Where the factor vanishes the
original expression is $+\infty$, which cannot lower an infimum. All
other original section weights are at least $\ell\min_K f$, because
canonical section norms are at most one.

The finite-dimensional filtered-space trace inequality is obtained by
extending a basis adapted to the subspace filtration to an adapted basis
of the whole space; its quotient weights have that lower bound. With
$r_\ell=h^0(X,\ell L)$ and $d_\ell=h^0(X_m,\ell A_m)$ it gives, up to
rounding errors bounded by the two ranks,
\[
 \operatorname{Tr}_{\ell,f}\geq
 \operatorname{Tr}_{\ell,\psi}
             +(r_\ell-d_\ell)\ell\min_K f.         \tag{8.8}
\]
First let $\ell$ tend to infinity with $m,f,\psi$ fixed. Section growth
and the bounded nef trace just identified give
\[
 \mathscr E(f)\geq e_m(\psi)+(\min_K f)(V-V_m).      \tag{8.9}
\]
The term on the right is generally negative. It cannot be discarded at
finite $m$, and its coefficient may not be replaced by a normalized
probability rank. Only the subsequent limit $m\to\infty$, for a fixed
finite continuous obstacle $f$, removes it. No integration of a moving
unbounded original section function is used in this argument.

\section{Quadratic relative contraction with its generic error}
On the auxiliary actual nef bundle $A_m$ consider a bounded model
$\psi\leq0$. Denote its mixed raw products by
$\mu_k=(A_m+dd^c\psi)^k\wedge A_m^{n-k}$, including every $b_E$.
Their masses are $V_m$. Set $a_k=\int(-\psi)\,d\mu_k\geq0$. The trivial
metric has measure $V_m\delta_{v_0}$, so
$a_0=-\psi(v_0)V_m$. Its energy cost is
\[
 B_m(\psi)=-e_m(\psi)=\frac{1}{n+1}\sum_{k=0}^n a_k.  \tag{9.1}
\]
For $0\leq r\leq1$, multilinearity on this fixed model gives the exact
finite polynomial
\[
 B_m(r\psi)=\frac r{n+1}\sum_{j=0}^n\sum_{k=0}^j
             {j\choose k}r^k(1-r)^{j-k}a_k
 \leq r a_0+C_n r^2B_m(\psi),\quad C_n=2^{n+1}.     \tag{9.2}
\]
The $k=0$ terms give at most $ra_0$; every remaining term contains another
factor $r$, and the sum of its binomial coefficients is bounded by
$C_n$. Thus (9.2) is proved for every dimension, rather than inferred from
a finite-dimensional experiment. The linear term vanishes only if the
generic value is exactly zero.

Let $u=P_\theta(f_D)$ be normalized and let
$u^{(r)}=(1-r)V_\theta+ru$. The actual section formulas show
\[
 (1-r)\phi_{m,0}+r\phi_{m,D}=q_{T_m}+r\psi_m\leq u^{(r)}. 
\]
Apply (8.9) to each \emph{fixed} continuous majorant $f$ of $u^{(r)}$,
with $\psi=r\psi_m$. Equations (8.6), (9.2) give
\[
 \mathscr E(f)\geq
   -r(-g_mV_m)-C_nr^2B_m(\psi_m)+(\min_K f)(V-V_m).
                                                               \tag{9.3}
\]
Now let $m$ tend along the complete divisibility subnet. The generic
error and the missing-rank error both tend to zero. Finally take the
infimum over continuous majorants. This proves
\[
 B(u^{(r)})\leq C_nr^2B(u),\qquad B(u)=-E(u).        \tag{9.4}
\]
For every normalized genuine $u\in\E^1$, its normalized model envelopes
$u_i=P_\theta(f_i)$ decrease all-point to $u$. One may choose $f_i\leq0$:
the original normalized $u$ is at most zero, and taking a minimum with
zero preserves a model majorant. Their generic values are zero because
they dominate $u$ and lie below $f_i\leq0$. Decreasing energy continuity
passes (9.4) to $u$. Rational $r$ suffices for the actual section
calculation; finite concavity and continuity on a segment give every real
$r\in[0,1]$. This argument retains the intrinsic singular reference
$V_\theta$ throughout.

\section{Intrinsic relative targets: integrability and coercivity}
Write $Z_\theta=\{V_\theta=-\infty\}$. For this section the target datum is
a positive original Radon measure $\nu$ of raw mass $V$ such that
$\nu(Z_\theta)=0$. For normalized $u\in\E^1$ define
\[
 I_\nu(u)=\int(V_\theta-u)\,d\nu\in[0,+\infty],\qquad
 D_\nu=\sup_{u(v_0)=0}\{I_\nu(u)-B(u)\}.             \tag{10.1}
\]
The relative difference is taken on the complement of $Z_\theta$ and
uses the actual values. No absolute integrability of $V_\theta$ is assumed.
The definition is constant invariant when written as
$E(u)-\int(u-V_\theta)d\nu$ for unnormalized potentials.

\begin{proposition}
In this intrinsic relative convention, $D_\nu<\infty$ if and only if
every actual $\E^1$ relative difference is $\nu$-integrable. A finite dual
target satisfies, for $0<\epsilon\leq C_n$,
\[
 I_\nu(u)\leq\epsilon B(u)+C_nD_\nu/\epsilon,\qquad
 I_\nu(u)-B(u)\leq-(1-\epsilon)B(u)+C_nD_\nu/\epsilon. \tag{10.2}
\]
\end{proposition}
\begin{proof}
Finite $D_\nu$ implies integrability directly from (10.1). Conversely,
suppose every relative difference is integrable. Their integrals are
uniformly bounded on $B(u)\leq1$. Otherwise choose normalized $u_j$ of
cost at most one with $I_\nu(u_j)\geq4^j$. The genuine functions
\[
 v_N=\sum_{j=1}^N2^{-j}u_j+2^{-N}V_\theta
\]
decrease, have generic value zero, and have cost at most one by concavity.
Sections 4 and 6 make their limit an actual $\E^1$ function. On the
complement of $Z_\theta$, its relative gap is the increasing sum
$\sum2^{-j}(V_\theta-u_j)$. Countable monotone convergence makes its
integral infinite, a contradiction. Denote the uniform bound by $C_\nu$.

For $B(u)\geq1/C_n$ apply (9.4) with
$r=(C_nB(u))^{-1/2}$. The contracted function has cost at most one and
$rI_\nu(u)\leq C_\nu$. For smaller costs the same unit-cost bound applies
directly. Thus $I_\nu(u)\leq C_\nu(1+\sqrt{C_nB(u)})$, whose difference
from $B(u)$ has finite supremum. This proves the converse.

Finally $I_\nu(u^{(r)})=rI_\nu(u)$ and (10.1), (9.4) imply
$rI_\nu(u)\leq C_nr^2B(u)+D_\nu$. Choose $r=\epsilon/C_n$ to obtain
(10.2). The reference gives $D_\nu\geq0$. Choosing $\epsilon<1$ shows
that any maximizing sequence has bounded actual energy cost. This last
bound alone does not assert that the objective attains its supremum.
\end{proof}

\section{A target-independent outer-derivative measure}
This section constructs an original-carrier positive measure from the
section-volume energy. Its identification with all unbounded weighted
Monge--Amp\`ere limits is a further obligation, explicitly left in
Section 13. The construction does not solve an inverse target problem.

\subsection{Deficient measures and their positive tail}
For a normalized nonpositive model obstacle $f_D$ with
$\mathscr E(f_D)\geq-B$, set
\[
 g_D(c)=\vol(\alpha_D(c))/(n+1),\quad
 \mu_{D,c}=\sum_E b_E\langle\alpha_D(c)^n\rangle\cdot E\,\delta_{x_E}.
                                                               \tag{11.1}
\]
At $c>B/V$, $\alpha_D(c)$ is big. Indeed the section formula gives
$g_D(c)\geq cV-B>0$. The same formula shows that $g_D$ is convex,
$V$-Lipschitz, $g_D(0)=0$, and eventually
$g_D(c)=cV+\mathscr E(f_D)$. Its derivative is the mass of $\mu_{D,c}$,
so
\[
 0\leq V-\mu_{D,c}(K)\leq B/c.                     \tag{11.2}
\]
Each coefficient in (11.1) is nonnegative: increasing the effective
divisor $E$ cannot decrease volume, and BFJ's differentiability computes
that derivative.

There is a coefficientwise comparison, stronger than (11.2). In the
lattice description the quantity
\[
 T_D(c)=\mathscr E(f_D)+cV-g_D(c)
       =\lim_m\frac{n!}{m^{n+1}}\sum_j\min(w_{mj}+mc,0)          \tag{11.3}
\]
is monotone when an effective vertical divisor is added to $D$.
This follows from inclusion of the actual lattices, or from their
finite-dimensional filtrations before the limit. Differentiate (11.3)
in each effective $E$ at a big $\alpha_D(c)$, using Section 6 for the
tail derivative. It follows that
\[
 0\leq\mu_{D,c}\leq\mu_D,
           \|\mu_D-\mu_{D,c}\|_{\rm TV}\leq B/c.     \tag{11.4}
\]
This is not an assertion that arbitrary positive products are linear.
Both measures refer to the same finite model, and the tail is identified
through its actual section weights.

\subsection{Uniform variation on a fixed test determination}
Fix a rational model direction $h$, represented by a signed vertical
Cartier $H$ on a smooth determination $Z$. Its value at $x_E$ is
the coefficient of the pulled $H$ divided by $b_E$. Refinements carrying
$D$ are allowed. For $D\leq0$ put $\gamma=\rho_X^*L+cF$, so
$\alpha_D(c)\leq\gamma$. Choose an ample $T$ on $Z$; its pullback to a
refinement is nef. Write
$M=(n+1)\langle\gamma^n\rangle\cdot T$ and
$\kappa=(n+1)(cV-B)>0$. Along $\alpha-rT$ while it is big,
volume differentiation and positive-product order bound the loss by
$rM$. Consequently
\[
 \alpha-r_0T\text{ is big for }r_0=\kappa/[2(M+1)].  \tag{11.5}
\]
If the path left the big cone earlier, volume continuity at that first
exit would contradict the positive lower bound $\kappa-rM$.
This argument is uniform over refinements and over these obstacles.

Choose $a$ with $aT\pm H$ nef on $Z$. Equation (11.5) makes
$\lambda\alpha\pm H$ pseudoeffective for $\lambda=a/r_0$.
For $|t|\lambda<1/2$ positive-product order and homogeneity sandwich
$\langle(\alpha+tH)^n\rangle$ between
$(1-\lambda|t|)^n\langle\alpha^n\rangle$ and
$(1+\lambda|t|)^n\langle\alpha^n\rangle$.
Pair this sandwich with the two fixed nef terms of a decomposition
$H=H_1-H_2$, never with $H$ directly. Their upper bounds use
$\langle\gamma^n\rangle$. Integrating the resulting derivative estimate
gives one finite $K_c$ such that
\[
 |g_{D+tH}(c)-g_D(c)-t\!\int h\,d\mu_{D,c}|
                                      \leq K_ct^2/2.             \tag{11.6}
\]
The constants depend on $B,c,L,h,Z$, but not on the obstacle or its
refinement. All uses of BFJ are in big algebraic classes on smooth
projective determinations.

Let $H_0=\|h\|_\infty$. Along either sign of $t$, subtract
$|t|H_0F$ from $D+tH$. Its obstacle is nonpositive and its energy is at
least $-B'$, where $B'=B+2VH_0$ when $|t|\leq1$.
It has deficient parameter $c+|t|H_0$. Apply (11.4) there and at $D$,
integrate the deficient derivative along the short path, and use (11.6).
One obtains
\[
 |\mathscr E(f_D+th)-\mathscr E(f_D)-t\!\int h\,d\mu_D|
 \leq K_ct^2/2+|t|H_0(B+B')/c.                     \tag{11.7}
\]
The change of time parameter restores exactly the constants subtracted
from the obstacle; its derivative is the deficient mass. In particular
its additional error is bounded by the same mass deficit in (11.4).
The order of limits in (11.7) is fixed $c$, then $t\to0$, then $c\to\infty$.
An estimate stated only relative to a polar reference would not establish
this original continuous-direction estimate.

\subsection{Entire actual bases and absolute continuous directions}
Let normalized $u\in\E^1$ have cost at most $B$. Take all nonpositive
model majorants $f_i\downarrow u$; their envelopes decrease to $u$ and
their energies decrease to $E(u)$. Their positive measures $\mu_{f_i}$
have mass $V$. Compactness of positive mass-$V$ Radon measures supplies
cluster measures on the \emph{original} compact carrier. For each fixed
model $h$,
\[
 P_\theta(f_i+th)\downarrow P_\theta(u+th).
\]
These functions lie between $u\pm|t|\|h\|_\infty$, hence are genuinely
finite energy. Pass (11.7) to any cluster measure and to their decreasing
energies. Divide by $t$ with each sign, then use the stated limit order.
This proves
\[
 \left.\frac d{dt}\right|_{0}E(P_\theta(u+th))=\int h\,dM(u).    \tag{11.8}
\]
All cluster measures have these same integrals. Model functions are
uniformly dense and masses are $V$, so the cluster is unique and the
whole net converges. This defines $M(u)$ without a target. If $h\in C(K)$,
uniform model approximation changes the divided energy increment by at
most $V\|h-h_j\|_\infty$, and changes the measure integral by the same
bound. Let $t\to0$ for fixed $j$, then $j\to\infty$. Both derivative
signs in (11.8) hold for \emph{absolute original continuous tests}.
Adding a constant to $u$ preserves $M(u)$ and changes energy by $cV$.

For ordered finite-energy endpoints $a\geq b$ of bounded cost, compare
their perturbed envelope energies using order. The same estimate gives
for a fixed model $h$
\[
 \left|\int h\,d(M(a)-M(b))\right|
 \leq \frac{E(a)-E(b)}{|t|}+K_ct+C_h(B)/c             \tag{11.9}
\]
for small positive $t$. It also follows by first comparing two model
obstacles in (11.7), then passing to their majorants. First send the
ordered energy gap to zero, then $t\to0$ for fixed $c$, then $c\to\infty$.
Uniform model density gives weak convergence against every original
continuous test. In particular $M$ is continuous along decreasing genuine
finite-energy nets with finite-energy endpoint. At a model envelope it
agrees with its finite atomic $\mu_D$: use model majorants $f_i$ of the
envelope and the obstacles $\min(f_D,f_i)$. Their energies equal
$E(P_\theta(f_D))$, by the envelope order identity, and (11.9) identifies
their measures with $\mu_D$ before taking their endpoint.

This is an actual outer-energy measure of full volume on the original
carrier. The term ``genuine full finite-energy Monge--Amp\`ere operator''
would additionally require the weighted and negligible-set identifications
listed below; they are not conclusions of weak convergence or (11.8).

\subsection{Algebraic polar sets: a limited nullness conclusion}
Fix a nonzero section $s$ of an actual effective divisor $S$ on $X$.
On the model of $D$, its pulled divisor has vertical coefficients
$b_Ex_E(s)$ and an effective horizontal remainder. For fixed $c>B/V$,
\[
 \int v(s)\,d\mu_{D,c}
 \leq\langle\alpha_D(c)^n\rangle\cdot\rho_X^*S
 \leq a\langle\gamma^n\rangle\cdot T=:C_c,           \tag{11.10}
\]
where $aT-\rho_X^*S$ is nef on a fixed determination. The first inequality
uses the nonnegative effective-horizontal volume derivative. Put
$b_N=\min(v(s),N)$, an original continuous nonnegative function, equal
to $N$ on the algebraic zero locus. Equation (11.4) gives
$\int b_Nd\mu_D\leq C_c+NB/c$. Passing to $M(u)$, dividing by $N$, and
letting $N\to\infty$ proves
$M(u)(S^{\an})\leq B/c$. Let $c\to\infty$ to obtain zero mass.

An actual Kodaira decomposition gives $V_\theta\geq q_E$ with effective
$E$; hence $Z_\theta\subset E^{\an}$ and $M(u)(Z_\theta)=0$.
This argument does \emph{not} give a bound on
$\int(V_\theta-u)dM(u)$, absolute reference integrability, nullness of
every nonalgebraic pluripolar set, or negligible discrepancy for arbitrary
suprema. Those conclusions cannot be substituted for (11.10).

\section{The exact additional zero-volume range and three safeguards}
\subsection{Actual numerically trivial rational bundles}
Let $n\geq1$ and let the \emph{actual} rational bundle $L$ be numerically
trivial, including a nontorsion element of $\operatorname{Pic}^0(X)$.
Use the same numerical varying-actual-bundle single-net cone with
$\theta=0$, and the canonical trivial metric as reference. The zero
function is genuine: for every positive rational $\epsilon$, the actual
bundle $L+\epsilon A$ is ample; its weight-zero generating systems have
identically zero scalar potential. Thus one constant scalar net with
classes tending to zero is a defining net. Constants follow by shifts.

Conversely normalize $u\in\PSH(0)$ by $u(v_0)=\sup u=0$. Integer
multiplication of each defining section system proves $mu\in\PSH(0)$
for every positive integer $m$. Tensoring into one fixed actual ample
bundle makes all $mu$ normalized genuine ample scalars. G0 compactness
gives a finite coordinate lower bound $-C_z$ at each fixed rational
divisorial $z$, independent of $m$. Thus
\[
 -C_z\leq mu(z)\leq0\quad\text{for every }m,
\]
which forces $u(z)=0$. G0 all-point order determination in that ample
bundle gives $u=0$ on the entire carrier. Consequently
\[
 \PSH(0)=\mathbf R.                                 \tag{12.1}
\]
This proof uses neither fixed-degree sections of $L$ nor decreasing
closure in a general pseudoeffective class. It matches the conclusion,
but replaces the decreasing-closure step, of BJ Corollary 4.24.

The actual bounded raw zero-class product is zero: its positive ample
perturbations have mass $\epsilon^n(A^n)\to0$; positivity bounds every
continuous test by this mass. The construction is target independent.
Its energy relative to zero is zero on constants, with zero constant
law. Hence the entire finite-energy space in this convention is
$\E^1(0)=\mathbf R$, the raw operator is $M_0(c)=0$, and the intrinsic
positive mass-zero target is necessarily the zero measure, of dual energy
zero. There is an actual bijection
\[
 \E^1(0)/\mathbf R=\{\text{one point}\}\longrightarrow\{0\}.
                                                               \tag{12.2}
\]
Solutions are unique modulo constants; comparison and both continuous
energy derivative signs are trivial. No mass is divided by zero. The
unnormalized zero energy does not define a separating metric on constants.
This is the \emph{whole numerically trivial rational-bundle range} on
the retained smooth projective $X$, not every zero-volume class.

For example on $X=\mathbf P^1\times\mathbf P^1$ let
$L=p_1^*\OO(1)$. It is nef, nonzero numerically, and $L^2=\vol(L)=0$.
The pullback of $\max(\log\|z_1\|_{\rm triv},-1)$ is a nonconstant bounded
generated-system potential; zero is another. Their bounded nef top
products are positive with zero mass and thus vanish. One can verify
this on each finite model by G0 ample perturbation and intersection
polynomials; the masses for $L+\epsilon\OO(1,1)$ are
$2\epsilon(1+\epsilon)\to0$. Consequently this raw zero equation cannot
have uniqueness modulo constants on its whole bounded cone. This example
is a limitation of the zero-mass top equation, not a disproof of the
requested positive-volume theory.

\subsection{A polar FS function need not have original finite energy}
On $\mathbf P^1$ with $\OO(1)$ let
$\phi=\log\|z_1\|_{\rm triv}$ and
$\phi_k=\max(\phi,-k)$. Each $\phi_k$ is a finite generating-system
potential, normalized at $v_0$. On the branch toward $\{z_1=0\}$ it
has slope $-1$ until the valuation $v_k$ with $v_k(z_1)=k$, and slope
zero thereafter. Its raw curve measure is $\delta_{v_k}$ of mass one.
The actual curve energy formula gives
\[
 E(\phi_k)=\tfrac12\left(\int\phi_kd\delta_{v_0}
                        +\int\phi_kd\delta_{v_k}\right)=-k/2.
\]
G0 decreasing energy continuity proves $E(\phi)=-\infty$. Thus the
physical single-section FS function is genuine psh but is outside the
original ample $\E^1$. In particular the polar complete-system functions
in (8.4) cannot enter original $\E^1$ estimates merely because they have
an algebraic moving decomposition.

\subsection{Volume is not self-intersection}
On the blowup of $\mathbf P^2$ at a point, let $H$ be the pulled line
class and $E$ the exceptional curve. The actual big bundle $L=2H+E$
has $L^2=4-1=3$. Its intersection with $E$ is $-1$. Every divisor in
$|mL|$ therefore contains $E$; repeating after subtraction forces
multiplicity at least $m$. Multiplication by $s_E^m$ gives the exact
complete-system identity
\[
 H^0(X,mL)=s_E^mH^0(X,2mH),\qquad
 h^0(X,mL)={2m+2\choose2}.
\]
It follows that $\vol(L)=4$. A raw measure of mass $3$ would have the
wrong section-volume constant law in this big nonnef example.

\subsection{The abstract twice-closure counterexample is preserved}
Let $K_0=\{0\}\cup\{1/r:r\geq1\}$ with its compact convergent-sequence
topology. For $i,j\geq1$ define an extended continuous function by
\[
 f_{ij}(0)=0,\quad f_{ij}(1/i)=-\infty,\quad
 f_{ij}(1/r)=-j\ (r\leq i,\ r\ne i),\quad
 f_{ij}(1/r)=0\ (r>i).
\]
Take the family consisting of these functions and all their constant
shifts, with a constant class label. For fixed $i$, its decreasing
$j$-limit $g_i$ is $-\infty$ on the first $i$ isolated points and zero
elsewhere. The $g_i$ decrease to $g$, which is $-\infty$ at every
isolated point and is zero at $0$.

A decreasing net of the original functions cannot contain two different
pole indices: a common subsequent function would need both poles, while
every available stage has just one. Thus any such net has one fixed pole
index. Its values at all later isolated points equal its value at $0$;
constant shifts preserve this equality. It cannot converge all-point to
$g$. The twice-decreasing closure is strictly larger than the single-net
closure. This is a mathematical counterexample to an abstract flattening
assertion, not to the algebraic relative-floor theorem of Section 4.
The latter uses its additional actual reference bound at old poles.

\section{Exact unproved steps for the unchanged original equation}
The following statements remain unproved \emph{in this paper}. They are
not mathematical grants from the supplied historical candidates.
\begin{enumerate}
\item \textbf{Unbounded genuine Monge--Amp\`ere identification.}
The measure $M(u)$ is a target-independent outer-energy measure with raw
volume and finite-model compatibility. Identification with every genuine
finite-energy weighted moving-product limit, including independence of
all simultaneous unbounded approximations, needs unbounded mixed moment
estimates. Positivity of finite nef models and weak top convergence alone
do not supply those estimates.
\item \textbf{Entire output integration and negligible discrepancy.}
Prove that $\int(V_\theta-v)dM(u)$ is finite for every normalized genuine
$v\in\E^1$, and that nonalgebraic polar and negligible discrepancies have
zero output mass. Algebraic section-zero nullness in Section 11 does not
prove these assertions. Absolute integrability of the reference, when
needed by a particular proposed identification, requires its own bound.
\item \textbf{Absolute continuous-test approximation.}
For a bounded-cost family of model outputs converging weakly to a given
intrinsic target, construct fixed original continuous approximants to
each unbounded relative test with controlled \emph{absolute} $L^1$ error
at both the family and the limit. Keep the polar-collar parameter fixed
before section degree and measure index. A signed relative energy bound
or uniform control relative to a singular reference is not this statement.
\item \textbf{Strong fixed-target integration.}
Show continuity of $u\mapsto\int(V_\theta-u)d\nu$ on normalized
energy-bounded sets in the genuine weak topology, for every target in
(10.1). Sections 9--10 prove integrability and coercivity, not continuity.
The weakly usc energy does not repair the wrong semicontinuity sign of
the integral in a maximizing sequence.
\item \textbf{Actual arbitrary-target attainment and the equation.}
Coercivity gives bounded cost and compactness gives an actual cluster
function. To prove that the cluster maximizes the objective, rule out
an energy/integral jump using the preceding integration theorem or an
independently justified genuine strong-completeness argument. Only then
may both signs of (11.8) identify its measure with the target. No target
may be used to define $M$ itself.
\item \textbf{Full comparison and every-solution uniqueness.}
Extend the finite atomic supporting planes to all $\E^1$ bases with all
relative integrals identified. Prove strict separation of the resulting
energy or $I$ quantities and original Borel strict-set mass comparison.
The strict set of two usc functions need not be open. Passing weak
measures directly through its indicator is invalid. Uniqueness of a
variational maximizer, even if proved, would not yet give uniqueness of
every operator solution without the entire supporting-plane theorem.
\end{enumerate}

The supplied components 143--171 and candidate193 propose mechanisms for
these items, including auxiliary ample H\"older estimates, capped section
tests, model maximizing sequences, and nonlinear Lipschitz tests. Their
full mathematical text is preserved. We have reconstructed Sections
3--11 directly, rather than treating the component numbers, historical
verdicts or hashes as premises. This paper has not independently supplied
the six items above. In particular BJ Lemma 7.30, as read, assumes an
actual ample class; its regularization and Cauchy--Schwarz proof is not
an all-E1 nonnef theorem after the word ``ample'' is removed. Its printed
telescoping sum starts at $j=1$; the actual algebraic difference of two
$n$th powers starts at $j=0$. This correction supplies one elementary
identity, not the missing nonample passage.

No complete original P16.5 equation theorem is claimed. The additional
zero-volume theorem is precisely (12.1)--(12.2), and there is no claim
for general nonzero nonbig pseudoeffective classes. All seven contextual
family items and the exact original completion criteria remain unchanged
in the frozen input. The canonical result is \texttt{candidate=false}.

\appendix
\section{Source applications, literal coordinates and read scope}
The authoritative source record is \texttt{source-ledger.json}; raw
connector responses and original renderer coordinates are in
\texttt{sources/}. The extracted excerpt files give their own one-based
LF coordinates and a map to the original labels. A renderer label
\texttt{L147}, or \texttt{L3615@P50}, is \emph{not} represented as LF148
of a locally downloaded PDF. The failed inherited-proxy PDF downloads
are retained separately. No such downloaded PDF or its byte hash is
invented.

G0 and G1 are the only caller-granted complete mathematical bodies.
Both were read in the supplied assignment and preserved as whole units.
G0 is used only through its actual ample foundations or a fixed-model
finite-polynomial ample perturbation. G1 is used only in the resolved
equivariant signature of Section 3, with full lattices, cofinal nonempty
degrees, supported poles and the actual generation twist. No general
nonample closure is read into either grant.

The pertinent primary units read for the new algebraic applications are:
BFJ v2, Section 2.2 (positive products, directed least upper bounds and
compactness), Section 2.3 (continuity and nef agreement), Proposition 2.13
(nef competitors with psef remainder), and the entire pertinent proof of
Theorem A in Section 3.2 (Morse estimate, corollary and both derivative
signs). Their underlying ordinary Hodge-index, resolution, Fujita,
Kodaira decomposition and section-growth theorems remain explicitly
imported algebraic characteristic-zero background. We do not claim a
recursive reading or reproof of those complete theories. BFJ's stated
field extension is at original renderer L147. Its HTML rendering prints
an automatically generated 2026 date; the version-pinned arXiv identifier
is math/0608260v2, posted 4 September 2007, and is the version used here.

For the ample nonnegligibility application we read BJ v4's actual model
calculation, Lemmas 2.28--2.29, the complete finite-vertex Lemma 4.26,
and Lemma 5.7 with the complete Theorem 5.6 proof. Their ample
regularization step is supplied by G0. No unrestricted use of the
abstract Lemma 4.6 is made. BJ Corollary 4.24 is recorded as a matched
zero-class conclusion, with the different proof in Section 12. The
complete pertinent Lemma 7.30 proof was read for gap diagnosis; its
ample domain is retained. The read capacity/negligibility passages do
not constitute a proved transfer to $M(u)$.

The initial targeted literature search found relevant complex and toric
results, but no theorem was imported from a search snippet. This is not
an exhaustive bibliographical nonexistence claim. The published BFJ
source, G0, G1, ordinary algebraic background and deductions in this
paper are distinguished in the source-application register.

\section{Artifacts and reproducible checks}
The exact original JSON, decoded task and decoded assignment are in
\texttt{frozen/}. Their extraction is reproduced by
\texttt{python3 freeze\_input.py}. The whole supplied mathematical units
and their inclusive assignment-LF coordinates are in
\texttt{supplied/index.json}. Run \texttt{python3 rerun.py} to check
preservation, source-unit maps, checksums, result schema and explicit
finite arithmetic, then compile this paper twice with pdfLaTeX.
Commands, environment versions, stdout/stderr, compilation logs and
artifact hashes are retained in \texttt{logs/} and
\texttt{artifact-manifest.json}. These checks establish reproducibility
and catch transcription errors; they do not certify the mathematics.

No agents, other campaigns, credentials, runtime configuration or
authentication were used or changed. All generated work is in the fresh
artifact directory. The cloud-runtime skill was read for read-only
network/environment checks. Its environment-status tool was unavailable;
the readable network policy and inherited-proxy failure are recorded.
No resource-heavy calculation was launched. No operator-provided
supervised compute service was exposed, so no such facility is asserted.

\begin{thebibliography}{9}
\bibitem{BJ} S. Boucksom and M. Jonsson,
\emph{Global pluripotential theory over a trivially valued field},
\href{https://arxiv.org/pdf/1801.08229v4}{arXiv:1801.08229v4} (2022).
The actual ample applications and diagnostic units are recorded above.
\bibitem{BFJ} S. Boucksom, C. Favre and M. Jonsson,
\emph{Differentiability of volumes of divisors and a problem of Teissier},
\href{https://arxiv.org/html/math/0608260v2}{arXiv:math/0608260v2} (2007).
Only ordinary algebraic positive products and volume differentiation are
imported; no nonample analytic equation theorem is contained in this use.
\bibitem{ADD} S. Boucksom and M. Jonsson,
\emph{Addendum to global pluripotential theory over a trivially valued field},
\href{https://arxiv.org/pdf/2206.07183v2}{arXiv:2206.07183v2}.
The actual lattice calculation is within G1; its general closure theorem
is not an additional premise here.
\bibitem{DEL} J.-P. Demailly, L. Ein and R. Lazarsfeld,
\emph{A subadditivity property of multiplier ideals},
\href{https://arxiv.org/html/math/0002035v2}{arXiv:math/0002035v2}.
The exact relative arbitrary-field application is granted only through G1.
\end{thebibliography}
\end{document}
